\documentclass[11pt,a4paper]{report}

\usepackage[T1]{fontenc}
\usepackage[utf8]{inputenc}
\usepackage[french]{babel}
\usepackage{lmodern}
\usepackage[a4paper,margin=2.2cm]{geometry}
\usepackage{xcolor}
\usepackage{titlesec}
\usepackage{tabularx}
\usepackage{booktabs}
\usepackage{array}
\usepackage{enumitem}
\usepackage{fancyhdr}
\usepackage{longtable}
\usepackage[hidelinks]{hyperref}

% ---- Couleurs de marque ----
\definecolor{sagiteal}{HTML}{00745A}
\definecolor{sagidark}{HTML}{1A1A2E}
\definecolor{sagigray}{HTML}{6B7280}
\definecolor{sagibg}{HTML}{F4F6F8}

\hypersetup{
  pdftitle={SAGI SCHOOL -- Documentation des modules},
  pdfauthor={HADY GESMAN},
  colorlinks=false,
}

% ---- Titres ----
\titleformat{\chapter}[hang]
  {\bfseries\Large\color{sagiteal}}{\thechapter.}{0.6em}{}
\titlespacing{\chapter}{0pt}{0pt}{12pt}
\titleformat{\section}
  {\bfseries\large\color{sagidark}}{\thesection}{0.6em}{}
\titleformat{\subsection}
  {\bfseries\normalsize\color{sagidark}}{\thesubsection}{0.6em}{}

% ---- En-tête / pied ----
\pagestyle{fancy}
\fancyhf{}
\renewcommand{\headrulewidth}{0.4pt}
\renewcommand{\footrulewidth}{0.4pt}
\fancyhead[L]{\small\color{sagiteal}\textbf{SAGI SCHOOL}}
\fancyhead[R]{\small\color{sagigray}Documentation des modules}
\fancyfoot[L]{\small\color{sagigray}by HADY GESMAN}
\fancyfoot[C]{\small\color{sagigray}Conforme SYSCOHADA Révisé}
\fancyfoot[R]{\small\thepage}

\newcommand{\feat}[1]{\item #1}
\newcommand{\modtag}[1]{{\small\color{sagigray}\textit{Route applicative : \texttt{#1}}}}
\setlist[itemize]{leftmargin=1.3em,itemsep=2pt,topsep=3pt}

\renewcommand{\tabularxcolumn}[1]{>{\small}m{#1}}

\begin{document}

% =========================== PAGE DE TITRE ===========================
\begin{titlepage}
\centering
\vspace*{2.5cm}
{\Huge\bfseries\color{sagiteal} SAGI SCHOOL\par}
\vspace{0.4cm}
{\Large\color{sagidark} Système de Gestion Scolaire\par}
\vspace{0.2cm}
{\large\color{sagigray} \textit{by HADY GESMAN}\par}
\vspace{2.2cm}
{\LARGE\bfseries\color{sagidark} Documentation des modules\par}
\vspace{0.5cm}
{\large\color{sagigray} Description fonctionnelle détaillée du système\par}
\vspace{3cm}

\begin{tabular}{rl}
\textbf{Version} & 1.11.2\\[4pt]
\textbf{Architecture} & Electron + Django + Angular\\[4pt]
\textbf{Déploiement} & Local (.exe / .deb) et Cloud (navigateur)\\[4pt]
\textbf{Conformité} & SYSCOHADA Révisé (OHADA)\\[4pt]
\textbf{Cadre RH/fiscal} & Réglementation du Sénégal\\
\end{tabular}

\vfill
{\small\color{sagigray} Document de référence --- édité le \today\par}
\end{titlepage}

\tableofcontents
\clearpage

% =========================== PRÉSENTATION ===========================
\chapter{Présentation générale}

SAGI SCHOOL est une application de gestion scolaire complète destinée aux
établissements d'enseignement. Elle couvre l'ensemble du cycle de gestion :
inscription des élèves, encaissement des frais de scolarité, comptabilité
conforme au référentiel \textbf{SYSCOHADA Révisé}, suivi académique
(notes et bulletins) et gestion de la paie du personnel selon la
réglementation sénégalaise.

\section{Architecture technique}
\begin{itemize}
\feat{\textbf{Bureau (local)} : application \textbf{Electron} embarquant un
serveur \textbf{Django} (Python) et une interface \textbf{Angular}. Distribuée
en installeur Windows (\texttt{.exe}) et Linux (\texttt{.deb} / AppImage).
Fonctionne \textbf{hors-ligne}.}
\feat{\textbf{Cloud} : la même interface Angular servie dans un navigateur,
adossée au backend Django hébergé. Les deux modes sont mutuellement exclusifs.}
\feat{\textbf{Multi-établissement (multi-tenant)} : chaque école est isolée
logiquement ; les données d'une école ne sont jamais visibles par une autre.}
\feat{\textbf{Base de données} : PostgreSQL.}
\end{itemize}

\section{Rôles et droits d'accès}
Le système distingue plusieurs profils utilisateur, du super-administrateur
éditeur jusqu'au simple lecteur :

\begin{center}
\begin{tabularx}{\textwidth}{>{\bfseries}l X}
\toprule
Rôle & Description \\
\midrule
Super Administrateur & Éditeur (HADY GESMAN) : gère l'ensemble des écoles clientes et leurs licences. \\
Administrateur École & Directeur d'une école cliente : accès complet aux modules de son établissement. \\
Responsable RH & Accès au module Ressources Humaines / Paie. \\
Responsable Comptable & Accès à la Comptabilité et au Fiscal. \\
Responsable Scolarité & Accès aux Élèves, Paiements et Académique. \\
Lecteur & Consultation seule. \\
\bottomrule
\end{tabularx}
\end{center}

\section{Licences et périmètre fonctionnel}
L'accès aux modules dépend du type de licence de l'établissement. Cinq formules
existent, dont une dédiée au programme \textbf{Taxawu Daara} (écoles coraniques) :

\begin{center}
\begin{tabularx}{\textwidth}{>{\bfseries}l X}
\toprule
Licence & Périmètre \\
\midrule
Essai (30 jours) & Découverte : tableau de bord, élèves, paiements, comptabilité, suivi mensuel. \\
Basic & Gestion essentielle : tableau de bord, élèves, paiements, suivi mensuel. \\
Pro & Basic + comptabilité SYSCOHADA complète. \\
Avancé & Périmètre complet : tout le Pro + académique, ressources humaines et fiscal. \\
Taxawu Daara & Offre dédiée Daaras : tableau de bord, élèves, paiements, comptabilité, académique, suivi mensuel. \\
\bottomrule
\end{tabularx}
\end{center}

\noindent Les modules \textbf{Ma Licence} et \textbf{Paramètres} restent
accessibles quelle que soit la licence. La matrice complète figure au
chapitre~\ref{ch:matrice}.

% =========================== MODULES ===========================
\chapter{Modules fonctionnels}

% ---------- Tableau de bord ----------
\section{Tableau de bord}
\modtag{/dashboard}

\noindent Vue de pilotage synthétique de l'établissement, affichée à la
connexion.
\begin{itemize}
\feat{\textbf{Indicateurs clés} : effectifs, recettes, ventilation des
encaissements et état des règlements.}
\feat{\textbf{Élèves à relancer} : liste des élèves présentant des arriérés,
détaillés \textbf{mois par mois}, avec reliquats et niveaux d'alerte
(légende et codes couleur unifiés avec le module Élèves).}
\feat{\textbf{Trésorerie} : suivi par canal (caisse, banque, mobile money).}
\feat{\textbf{Charges et marge} : aperçu des charges et du résultat
intermédiaire.}
\feat{\textbf{Journal d'audit} : traçabilité des actions sensibles
(qui a fait quoi, et quand).}
\end{itemize}

% ---------- Élèves ----------
\section{Élèves}
\modtag{/eleves}

\noindent Gestion du fichier élèves et de leur tarification.
\begin{itemize}
\feat{\textbf{Inscriptions} : matricule (numérotation automatique), identité
(nom, genre, date et lieu de naissance), \textbf{coordonnées des parents}
(père et mère avec téléphones), date d'inscription et statut.}
\feat{\textbf{Sections} : regroupements tarifaires portant les frais
d'inscription, de mensualité, d'uniforme et de fournitures.}
\feat{\textbf{Services optionnels} : cantine, transport, etc., avec
périodicité (mensuelle\ldots) et \textbf{abonnement par élève}.}
\feat{\textbf{Frais flexibles} : mensualité $\times$ nombre de mensualités,
avec \textbf{prorata} selon la date d'entrée et les mois concernés.}
\feat{\textbf{Prise en charge (PEC)} : gestion des bourses et sponsors via des
\textbf{taux} distincts sur l'inscription et la mensualité, avec observations.}
\feat{\textbf{Éditions PDF} : fiche élève (et édition), situation détaillée de
l'élève, \textbf{certificat de scolarité} et liste des élèves.}
\end{itemize}

% ---------- Paiements ----------
\section{Paiements}
\modtag{/paiements}

\noindent Encaissement des frais et émission des reçus.
\begin{itemize}
\feat{\textbf{Exercices} : année scolaire, dates de début/fin, nombre de
mensualités, soldes initiaux (caisse, banque, mobile), devise et clôture.}
\feat{\textbf{Saisie des règlements} : numéro de pièce unique et ventilation
fine (inscription, mensualité, uniforme, fournitures, cantine, divers),
avec mémorisation des \textbf{mois réglés} et des \textbf{services réglés}.}
\feat{\textbf{Canaux de paiement} : espèces, Orange Money, Wave, Free Money,
virement\ldots}
\feat{\textbf{Reçus PDF} en trois formats : \textbf{A4}, \textbf{A5} et
\textbf{ticket 80\,mm} (imprimante thermique).}
\feat{\textbf{Annulations} : un règlement peut être annulé ; les
\textbf{écritures comptables de contre-passation} conformes SYSCOHADA sont
générées automatiquement.}
\feat{\textbf{Comptabilisation automatique} : chaque encaissement alimente le
journal comptable (lien source / pièce).}
\end{itemize}

% ---------- Suivi mensuel ----------
\section{Suivi mensuel}
\modtag{/suivi-mensuel}

\noindent Tableau de bord des règlements mois par mois.
\begin{itemize}
\feat{Vue récapitulative des paiements et \textbf{reliquats} par élève sur
l'ensemble des mensualités de l'exercice.}
\feat{Identification rapide des \textbf{arriérés} pour relance, en cohérence
avec le tableau de bord.}
\feat{Totaux et synthèse par période.}
\end{itemize}

% ---------- Comptabilité ----------
\section{Comptabilité (SYSCOHADA Révisé)}
\modtag{/comptabilite}

\noindent Comptabilité d'engagement complète, conforme au référentiel
SYSCOHADA Révisé.
\begin{itemize}
\feat{\textbf{Plan comptable} : comptes par classe (1 à 9), typés
(charge, produit, actif\ldots), avec des comptes \textbf{système} non
supprimables.}
\feat{\textbf{Journal} : écritures débit/crédit horodatées et numérotées,
issues de la saisie ou générées automatiquement par les autres modules.}
\feat{\textbf{Grand livre} et \textbf{Balance} générale.}
\feat{\textbf{États de synthèse} : \textbf{Bilan} (actif / passif),
\textbf{Compte de résultat}, \textbf{Soldes Intermédiaires de Gestion (SIG)}
et \textbf{Notes annexes}.}
\feat{\textbf{Budget} : prévisions par compte avec ventilation \textbf{mensuelle}
(12 mois) et distinction charges fixes / variables.}
\feat{\textbf{Immobilisations} : suivi des biens, \textbf{amortissement}
(linéaire), durée d'utilisation, comptes d'immobilisation et d'amortissement,
cumul, \textbf{cession}, et règlement au comptant ou à crédit.}
\feat{\textbf{Clôture d'exercice} et export PDF de tous les états.}
\end{itemize}

% ---------- Fiscal ----------
\section{Fiscal}
\modtag{/fiscal}

\noindent Suivi des obligations fiscales et sociales (Sénégal).
\begin{itemize}
\feat{États et déclarations relatifs aux prélèvements : \textbf{BRS},
\textbf{CFCE}, \textbf{CSS}, \textbf{IPRES}, \textbf{IR}\ldots}
\feat{Paramètres fiscaux servant de base aux calculs (taux, plafonds).}
\feat{Exports PDF dédiés.}
\end{itemize}

% ---------- Académique ----------
\section{Académique}
\modtag{/academique}

\noindent Gestion pédagogique : notes, évaluations et bulletins.
\begin{itemize}
\feat{\textbf{Niveaux et classes} : niveaux scolaires (note maximale
paramétrable) et classes rattachées, avec effectifs.}
\feat{\textbf{Matières} : coefficient, note maximale et statut actif/inactif.}
\feat{\textbf{Types d'évaluation} : devoir, composition, interrogation\ldots
avec \textbf{poids} respectifs.}
\feat{\textbf{Évaluations et notes} : saisie par matière et par période, avec
gestion des \textbf{absences} et observations.}
\feat{\textbf{Période configurable} : trimestre, semestre ou période libre
(le nombre de périodes est paramétrable par établissement).}
\feat{\textbf{Bulletins} : moyennes, points, \textbf{rang par matière} et
appréciations, avec \textbf{bulletin PDF} détaillé et historique.}
\feat{\textbf{Analyse académique} : taux de réussite et statistiques de classe.}
\end{itemize}

% ---------- RH ----------
\section{Ressources Humaines / Paie}
\modtag{/rh}

\noindent Gestion du personnel et de la paie selon la réglementation
sénégalaise.
\begin{itemize}
\feat{\textbf{Employés} : matricule, type (enseignant, administratif\ldots),
poste, type de contrat (CDI/CDD), date d'embauche, salaire de base, situation
matrimoniale, nombre d'enfants, statut cadre, mode de paiement et
autorisation d'enseigner.}
\feat{\textbf{Paie} : calcul du salaire brut, des cotisations
\textbf{IPRES}, \textbf{CSS}, de l'\textbf{IR} et des autres retenues, jusqu'au
salaire net, par mois.}
\feat{\textbf{Paramètres fiscaux} : taux IPRES (général et cadre, parts
salariale et patronale), plafonds, CSS, barème IR, CFCE et ATMP --- versionnés
par année.}
\feat{\textbf{Avances sur salaire} et régularisations.}
\feat{\textbf{Bulletins de paie PDF} conformes.}
\end{itemize}

% ---------- Licences & Clients ----------
\section{Licences \& Clients}
\modtag{/licences}

\noindent Module réservé au \textbf{Super Administrateur} (HADY GESMAN) pour
l'administration commerciale de la solution.
\begin{itemize}
\feat{Gestion des \textbf{écoles clientes} (création, édition, activation).}
\feat{Délivrance et suivi des \textbf{licences} : type, clé unique, dates de
début et de fin, statut (active, expirée, suspendue, essai).}
\feat{\textbf{Tarification} : tarifs mensuels et cycle de facturation
mensuel ou annuel (remise de 10\,\% sur l'annuel).}
\feat{Pilotage à distance du périmètre fonctionnel ouvert à chaque client.}
\end{itemize}

% ---------- Ma Licence ----------
\section{Ma Licence}
\modtag{/ma-licence}

\noindent Vue de la licence active de l'établissement (toujours accessible).
\begin{itemize}
\feat{Type de licence, plan et niveau de support.}
\feat{Date d'expiration et statut.}
\feat{Information de \textbf{reconnexion} en mode local (tolérance hors-ligne
de 30 jours).}
\end{itemize}

% ---------- Paramètres ----------
\section{Paramètres}
\modtag{/parametres}

\noindent Configuration de l'établissement (toujours accessible).
\begin{itemize}
\feat{\textbf{Identité} : nom, ville, adresse, \textbf{RCCM}, \textbf{NINEA},
téléphone, e-mail et code d'établissement.}
\feat{\textbf{Logo} : import (intégré aux documents PDF officiels --- reçus,
bulletins, certificats, états comptables).}
\feat{\textbf{Préférences} : régime de paie, type de période scolaire et
nombre de périodes.}
\end{itemize}

% =========================== TRANSVERSES ===========================
\chapter{Fonctions transverses}

\section{Documents et impressions}
\begin{itemize}
\feat{Génération de \textbf{PDF} pour l'ensemble des documents officiels :
reçus, certificats, bulletins (élève et paie), fiches et états comptables.}
\feat{\textbf{Logo de l'établissement} embarqué automatiquement dans chaque
document.}
\feat{Mise en page optimisée pour la \textbf{lisibilité à l'impression}
(police, contraste, mise en gras), y compris en noir et blanc.}
\end{itemize}

\section{Multilingue (i18n)}
\begin{itemize}
\feat{Interface disponible en \textbf{français}, \textbf{anglais} et
\textbf{arabe}.}
\feat{Traductions versionnées à chaque build pour une prise en compte
immédiate après mise à jour.}
\end{itemize}

\section{Sécurité et traçabilité}
\begin{itemize}
\feat{\textbf{Isolation multi-tenant} systématique des données par
établissement.}
\feat{\textbf{Journal d'audit} des opérations sensibles.}
\feat{Gestion des \textbf{rôles} et des droits d'accès par module.}
\feat{Licence locale avec contrôle périodique et tolérance hors-ligne.}
\end{itemize}

% =========================== MATRICE ===========================
\chapter{Matrice modules / licences}
\label{ch:matrice}

Le tableau ci-dessous récapitule l'ouverture des modules selon le type de
licence. \textbf{Oui} indique que le module est accessible.

\begin{center}
\renewcommand{\arraystretch}{1.3}
\begin{tabularx}{\textwidth}{>{\bfseries\small}l *{5}{>{\centering\arraybackslash\small}X}}
\toprule
Module & Essai & Basic & Pro & Avancé & Taxawu \\
\midrule
Tableau de bord & Oui & Oui & Oui & Oui & Oui \\
Élèves          & Oui & Oui & Oui & Oui & Oui \\
Paiements       & Oui & Oui & Oui & Oui & Oui \\
Suivi mensuel   & Oui & Oui & Oui & Oui & Oui \\
Comptabilité    & Oui & --- & Oui & Oui & Oui \\
Académique      & --- & --- & --- & Oui & Oui \\
Ressources Humaines & --- & --- & --- & Oui & --- \\
Fiscal          & --- & --- & --- & Oui & --- \\
\midrule
Ma Licence      & Oui & Oui & Oui & Oui & Oui \\
Paramètres      & Oui & Oui & Oui & Oui & Oui \\
\bottomrule
\end{tabularx}
\end{center}

\vspace{1cm}
\begin{center}
{\color{sagigray}\small SAGI SCHOOL --- by HADY GESMAN --- Conforme SYSCOHADA Révisé}
\end{center}

\end{document}
